\section{Síntesis y ampliación}

Esta sesión puede condensarse en una «cadena explicativa de curvas». Primero, la language-model loss mejora de forma suave con compute, data y parameters; segundo, la task metric, debido a la chance baseline, el exact match y los cuellos de botella de múltiples subhabilidades, puede permanecer plana durante mucho tiempo y después subir; tercero, este umbral se desplaza con la model family, la data quality, la architecture, el fine-tuning y el measurement axis; cuarto, CoT cambia la elicitation interface y permite que los modelos más grandes desplieguen cálculos compuestos mediante tokens intermedios; quinto, self-consistency agrega además las respuestas con un inference budget de múltiples trayectorias; sexto, todas las conclusiones deben informarse junto con el costo, la benchmark construction y la configuración histórica.

\subsection{Una tabla de términos clave}

Los siguientes términos suelen confundirse dentro de una misma oración. Solo al separarlos se puede determinar si un artículo aumentó los recursos de entrenamiento, la capacidad del modelo, la información del prompt o la búsqueda en tiempo de inferencia.

\begin{longtable}{L{0.19\textwidth}L{0.32\textwidth}L{0.40\textwidth}}
\toprule
\textbf{Término} & \textbf{Problema que resuelve} & \textbf{Mecanismo central y relación con el curso} \\
\midrule
\endhead
Scaling law & Predecir cómo cambia la loss promedio en función de los recursos & Ajusta, sobre una familia de modelos controlada, la relación suave entre compute, data, parameters y loss; es la línea base de comparación para la emergence. \\
Emergent ability & Describir capacidades en tareas que la metric no detecta en modelos pequeños pero sí en modelos grandes & Depende de la chance/human baseline, el prompt, la metric y la model family; no equivale a haber demostrado un módulo interno discreto. \\
Few-shot prompting & Indicar al modelo el formato de la tarea sin actualizar los pesos & Coloca exemplars en el context para que el next-token model induzca la correspondencia entre entradas y salidas. \\
Instruction fine-tuning & Lograr que el modelo siga tareas en lenguaje natural de manera más estable & Actualiza los weights; puede inducir anticipadamente una desired behavior ya conocida en modelos más pequeños. \\
RLHF & Lograr que las salidas se ajusten mejor a la preference humana & Usa una comparison/reward signal para actualizar el modelo; difiere de CoT, que solo cambia el context. \\
Chain-of-thought & Proporcionar estados intermedios generables para tareas de varios pasos & Escribe una reasoning trace en los ejemplos para que el modelo despliegue la decomposition en el orden de los tokens. \\
Perplexity & Medir la next-token prediction sobre datos held-out & $\exp(\text{cross-entropy})$; sirve como proxy válido de la calidad del modelo, pero no es una puntuación global de capacidad en tareas posteriores. \\
Self-consistency & Reducir el riesgo de que una sola CoT se desvíe por azar & Muestrea varias reasoning paths y vota sobre la final answer; intercambia inference cost por accuracy. \\
Temperature & Controlar la aleatoriedad del next-token sampling & Cambia el grado de suavizado del softmax de los logits y aporta path diversity a la self-consistency. \\
\bottomrule
\end{longtable}

\subsection{Cómo evaluar nuevas afirmaciones sobre Scaling con el marco de esta sesión}

Ante una nueva gráfica ``emergent'', primero hay que pedir a los autores originales que publiquen cada checkpoint, prompt, baseline, seed y confidence interval, y luego verificar si la metric convierte una calidad continua en un umbral rígido; ante un nuevo reasoning method, hay que exigir que se informen a la vez el base model, los generated tokens, el sample count, la latency y las external tools; ante conclusiones entre distintas familias de modelos, es obligatorio preguntar si también cambiaron los data y el post-training. Solo así se puede convertir lo que «parece repentino» en un hecho de ingeniería reproducible.

\begin{warningbox}{Las tres conclusiones excesivas más frecuentes}
Primero, inferir una phase transition precisa a partir de unos cuantos checkpoints dispersos; segundo, atribuir por completo a los parámetros un cambio que en realidad combina parámetros, datos, arquitectura y post-training; tercero, tomar una CoT legible como el pensamiento interno real del modelo. Las tres van más allá de la evidencia presentada en clase y, al escribir y al evaluar, deben moderarse de forma deliberada.
\end{warningbox}

\subsection{Lecturas adicionales}

Los siguientes materiales son fuentes primarias que ya eran públicas antes de la fecha de la clase y que constituyen directamente la cadena de evidencia de esta sesión. Al leerlos, se recomienda reproducir primero los axes, la baseline y el prompt de una curva, y después comparar qué tan fuerte es la formulación de los autores sobre el mecanismo.

\begin{itemize}
  \item \href{https://arxiv.org/abs/2001.08361}{Kaplan et al., Scaling Laws for Neural Language Models}: la referencia clásica de loss scaling suave.
  \item \href{https://arxiv.org/abs/2206.07682}{Wei et al., Emergent Abilities of Large Language Models}: catálogo de emergencia de tareas y de prompting-techniques.
  \item \href{https://arxiv.org/abs/2211.02011}{Wei, Tay, and Le, Inverse Scaling Can Become U-Shaped}: por qué una tendencia negativa local puede revertirse.
  \item \href{https://arxiv.org/abs/2201.11903}{Wei et al., Chain-of-Thought Prompting Elicits Reasoning in Large Language Models}: el artículo principal sobre CoT.
  \item \href{https://arxiv.org/abs/2210.09261}{Suzgun et al., Challenging BIG-Bench Tasks and Whether Chain-of-Thought Can Solve Them}: BBH, heterogeneidad de tareas y scale threshold.
  \item \href{https://arxiv.org/abs/2210.03057}{Shi et al., Language Models are Multilingual Chain-of-Thought Reasoners}: MGSM y composicionalidad entre idiomas.
  \item \href{https://arxiv.org/abs/2203.11171}{Wang et al., Self-Consistency Improves Chain of Thought Reasoning in Language Models}: muestreo de múltiples trayectorias y votación por mayoría.
  \item \href{https://arxiv.org/abs/2210.11416}{Chung et al., Scaling Instruction-Finetuned Language Models}: evidencia de postentrenamiento sobre el scaling de task, model y CoT data.
\end{itemize}

En última instancia, el aprendizaje más importante de esta sesión no es memorizar algún umbral de parámetros, sino adquirir una disciplina ante las curvas: primero preguntar qué se mide, luego qué se cambia y solo al final si es posible extrapolar. El Scaling lleva al modelo a nuevas regiones alcanzables, los datos y el postentrenamiento desplazan las fronteras, y CoT y self-consistency cambian la manera en que exploramos esas regiones; una conclusión verdaderamente confiable debe explicar a la vez la capacidad, la interfaz, el costo y la incertidumbre.

\subsection{Autoevaluación posterior a la clase}

Después de leer esta sesión, el lector debe poder desglosar de forma independiente una nueva afirmación sobre capacidades, sin apoyarse en el lema «bigger is better». Las siguientes preguntas pueden servir como lista de verificación mínima para la lectura de artículos, el diseño de experimentos y la evaluación de productos.

\begin{importantbox}{Seis preguntas de autoevaluación}
\begin{enumerate}
  \item ¿Por qué una mejora suave de la cross-entropy puede ocurrir al mismo tiempo que un aumento repentino de la exact-match accuracy?
  \item ¿Qué supuestos se mezclan, respectivamente, en el número de parámetros, los FLOPs de entrenamiento, la cantidad de datos y la perplexity?
  \item ¿Cómo se puede usar una controlled ablation para distinguir entre better data y un model-family confounder?
  \item ¿Cuál de estos elementos cambia CoT: los weights, el context o el generated-token budget?
  \item ¿Bajo qué estructura de correlación de errores falla la self-consistency, incluso con un sample count muy grande?
  \item ¿Qué baselines, scale points o cost information le faltan todavía a una gráfica de emergence para sustentar decisiones de ingeniería?
\end{enumerate}
\end{importantbox}

\end{document}
